\subsection{Ordered rings}

DEFINITION 1. --- Given a commutative ring A, we say that an ordering on A is compatible with the ring structure of A if it is compatible with the additive group structure of A, and if it satisfies the following axiom:

(OR) The relations \(x \geqslant 0\) and \(y \geqslant 0\) imply \(xy \geqslant 0\) .

The ring A, together with such an ordering, is called an ordered ring.

Examples. --- 1) The rings Q and Z, with the usual orderings, are ordered rings. 2) A product of ordered rings, equipped with the product ordering, is an ordered ring. In particular the ring \(A^{E}\) of mappings from a set E to an ordered ring A is an ordered ring.

\begin{enumerate}
\def\labelenumi{\arabic{enumi})}
\setcounter{enumi}{2}
\tightlist
\item
  A subring of an ordered ring, with the induced ordering, is an ordered ring.
\end{enumerate}

In an ordered ring, the relations \(x \geqslant y\) and \(z \geqslant 0\) imply \(xz \geqslant yz\) . Indeed these inequalities are equivalent to \(x - y \geqslant 0\) , \(z \geqslant 0\) and \((x - y)z \geqslant 0\) respectively.

Analogously we can show that the relations \(x \leqslant 0\) and \(y \geqslant 0\) (resp. \(y \leqslant 0\) ) imply \(xy \leqslant 0\) (resp. \(xy \geqslant 0\) ). These results are often invoked under the name of sign rules (two elements are said to have the same sign if they are both \(\geqslant 0\) or both \(\leqslant 0\) ). They imply that, if A is a totally ordered ring, then every square is positive, and in particular that every idempotent (for example the unit element) is positive.

Example. --- There is only one totally ordered ring structure on Z: indeed 1 \textgreater{} 0, whence n \textgreater{} 0 for every natural number \(n \neq 0\) , by induction. In contrast there exist ordered ring structures on Z which are not totally ordered (see below).

Let P be the set of positive elements of an ordered ring A. It is known (VI, p.~3, Prop. 3) that P determines the ordering on A. To say that A is an ordered ring is equivalent to saying that P satisfies the following properties:

\[
\begin{array}{l l} \left(\mathrm{AP} _ {\mathrm{I}}\right) & \mathrm{P} + \mathrm{P} \subset \mathrm{P} \\ \left(\mathrm{AP} _ {\mathrm{II}}\right) & \mathrm{PP} \subset \mathrm{P} \\ \left(\mathrm{AP} _ {\mathrm{III}}\right) & \mathrm{P} \cap (- \mathrm{P}) = \{0 \}. \end{array}
\]

Indeed \((\mathrm{AP}_{\mathrm{I}})\) and \((\mathrm{AP}_{\mathrm{III}})\) state that the additive group of A is an ordered group (VI, p.~3, Prop. 3), while \((\mathrm{AP}_{\mathrm{II}})\) is a translation of (OR).

Recall that the following condition is necessary and sufficient for the order relation on A to be total :

\[
\left(\mathrm{AP} _ {\mathrm{IV}}\right) \quad \mathrm{P} \cup (- \mathrm{P}) = \mathrm{A}.
\]

Example. --- In Z, if we take P to be the set of positive (in the usual sense) even integers, we get a ring which is not totally ordered.

Recall also that, in a totally ordered abelian group, the relation \(n \cdot x = 0\) (for a natural number \(n \neq 0\) ) implies \(x = 0\) (VI, p.~4); this gives us the following result.

PROPOSITION 1. --- A totally ordered ring is torsion free as a Z-module (II, p.~313).

